% =====================================================================
\section{Автономное питание: батареи и аккумуляторы}
% =====================================================================
\lessonintro
  {выбрать источник питания для устройства без USB, безопасно обращаться с литиевым аккумулятором и измерять его заряд}
  {урок 5 (делитель напряжения и АЦП), урок 9 (сон и расчёт времени работы), урок 16 (раздельное питание)}
  {устройство работает без компьютера, показывает напряжение батареи и засыпает при разряде}
  {держатель на три батарейки AA или powerbank; для опытов с Li-ion --- аккумулятор с платой защиты и модуль зарядки TP4056}

\subsection{Сначала: откуда брать энергию}

Пока плата подключена к компьютеру, о питании можно не думать. Но готовое устройство --- метеостанция
на балконе, робот, ночник --- должно работать само. Для этого нужно ответить на три вопроса: какой
источник выбрать, как получить из его напряжения нужные плате \SI{5}{V} или \SI{3.3}{V} и как узнать,
что он скоро разрядится.

\subsection{Источники питания}

\begin{center}\small
\begin{tabularx}{\textwidth}{@{}l l X X@{}}
\toprule
Источник & Напряжение & Плюсы & Минусы \\
\midrule
Powerbank & 5 В (USB) & безопасен, готовая зарядка & может отключаться при малом токе спящей платы \\
3 × AA (щелочные) & 4,5 В → 3,3 В & просты, безопасны & не перезаряжаются; напряжение падает по мере разряда \\
4 × AA (NiMH) & 4,8 В & перезаряжаются, безопасны & нужно отдельное зарядное устройство \\
Li-ion / Li-Po & 4,2 В → 3,0 В & большая ёмкость при малом размере & требует защиты и аккуратности \\
\bottomrule
\end{tabularx}
\end{center}

Для первых автономных опытов возьмите powerbank или три батарейки AA, подключённые к выводу 5V и GND:
на плате есть стабилизатор, который сделает из них \SI{3.3}{V}.

\begin{caution}
Литиевый аккумулятор при коротком замыкании, перезаряде или повреждении корпуса может загореться.
Используйте только аккумуляторы с платой защиты, заряжайте их только специальным модулем и только под
присмотром, не носите в кармане без чехла, не паяйте провода прямо к корпусу. Вздувшийся или помятый
аккумулятор использовать нельзя. В учебных опытах Li-ion применяется только с разрешения педагога.
\end{caution}

\subsection{Цепочка питания}

\begin{figure}[H]
\centering
\begin{tikzpicture}[font=\sffamily\scriptsize,node distance=9mm]
  \node[blk,fill=blkper,minimum width=22mm,minimum height=12mm] (usb) {USB 5 В\\ зарядка};
  \node[blk,fill=blkrtc,minimum width=26mm,minimum height=12mm,right=of usb] (chg) {TP4056\\ заряд + защита};
  \node[blk,fill=blkrtc,minimum width=24mm,minimum height=12mm,below=10mm of chg] (bat) {Li-ion\\ 3,0--4,2 В};
  \node[blk,fill=blkrtc,minimum width=28mm,minimum height=12mm,right=of chg] (reg) {Преобразователь\\ в 5 В или 3,3 В};
  \node[blk,fill=blkcpu,minimum width=24mm,minimum height=12mm,right=of reg] (esp) {ESP32-S3\\ вывод 5V или 3V3};
  \draw[arr] (usb) -- (chg);
  \draw[darr] (chg) -- (bat);
  \draw[arr] (chg) -- node[above]{OUT} (reg);
  \draw[arr] (reg) -- (esp);
  \draw[arr,accent] (bat.east) -| node[pos=0.25,below]{делитель → АЦП (GPIO2)} (esp.south);
\end{tikzpicture}
\caption{Питание от литиевого аккумулятора. Модуль TP4056 заряжает аккумулятор и защищает его от
перезаряда, глубокого разряда и короткого замыкания; преобразователь даёт плате стабильное напряжение}
\end{figure}

Напряжение аккумулятора меняется от \SI{4.2}{V} до \SI{3.0}{V}, а плате нужно ровно \SI{3.3}{V}.
Есть два пути. Повышающий преобразователь делает из напряжения аккумулятора \SI{5}{V}, которые подают
на вывод 5V платы. Либо стабилизатор с малым падением напряжения выдаёт \SI{3.3}{V} прямо на вывод 3V3;
стабилизатор на самой отладочной плате для этого обычно не подходит: ему нужен запас около вольта.

\begin{tip}
Отладочная плата не рассчитана на экономию: её светодиод питания и микросхема USB потребляют несколько
миллиампер даже тогда, когда чип спит. Для устройства, которое должно работать месяцами, берут плату
без лишних деталей или отдельный модуль ESP32-S3.
\end{tip}

\subsection{Сколько осталось заряда}

Напряжение литиевого аккумулятора падает по мере разряда, поэтому заряд можно оценить вольтметром
из урока 5. Напряжение до \SI{4.2}{V} выше предела АЦП, и его делят пополам двумя резисторами по
\SI{100}{k\ohm}: большие номиналы почти не разряжают аккумулятор.

\begin{figure}[H]
\centering
\begin{tikzpicture}
\begin{axis}[
  width=0.62\textwidth, height=5.2cm,
  xlabel={остаток заряда, \%}, ylabel={напряжение, В},
  xmin=0, xmax=100, ymin=3.0, ymax=4.3,
  x dir=reverse,
  grid=major, grid style={gray!25},
  tick label style={font=\scriptsize}, label style={font=\small},
]
  \addplot[very thick,esp,mark=*,mark size=1.6pt] coordinates {(100,4.20) (80,3.95) (55,3.80) (30,3.70) (10,3.60) (0,3.30)};
  \draw[dashed,accent] (axis cs:100,3.4) -- (axis cs:0,3.4) node[pos=0.3,above,font=\scriptsize,text=accent]{порог «пора выключаться»};
\end{axis}
\end{tikzpicture}
\caption{Примерная зависимость напряжения Li-ion аккумулятора от остатка заряда. В середине кривая
почти плоская, поэтому оценка в процентах всегда грубая}
\end{figure}

\codecap{Измерение напряжения батареи и защита от глубокого разряда}
\codeinput{l27_battery}

Программа оценивает проценты по таблице из шести точек, а между точками проводит прямые. Когда
напряжение падает ниже \SI{3.4}{V}, устройство сохраняет данные и уходит в глубокий сон (урок 9): так
оно не разрядит аккумулятор «в ноль».

\begin{caution}
Измеряйте батарею без нагрузки или всегда в одинаковых условиях. Во время передачи по Wi-Fi напряжение
проседает, и программа может решить, что аккумулятор пуст, хотя заряда ещё достаточно.
\end{caution}

\subsection{Сколько проработает устройство}

Время работы оценивают делением ёмкости на средний ток: аккумулятор \SI{2000}{mAh} при токе \SI{100}{mA}
проработает около 20 часов. Формула среднего тока для устройства, которое спит и просыпается, разобрана
в уроке 9. Реальное время на 20--30 \% меньше расчётного: часть энергии теряется в преобразователе,
а полностью разряжать аккумулятор нельзя.

\subsection{Самостоятельный опыт: индикатор заряда}

Запитайте плату от трёх батареек AA через вывод 5V и измеряйте их напряжение через тот же делитель.
Для батареек пороги другие: около \SI{4.5}{V} --- новые, ниже \SI{3.6}{V} --- пора менять. Покажите
заряд тремя светодиодами (или полоской на OLED) и убедитесь, что устройство работает без компьютера.

\begin{trouble}
\begin{itemize}[leftmargin=*]
  \item \textbf{От батареи плата перезагружается при включении Wi-Fi} --- источник не выдерживает пики
        тока; добавьте конденсатор \SIrange{470}{1000}{\micro F} у выводов питания или возьмите источник мощнее.
  \item \textbf{Powerbank отключается через полминуты} --- он считает малый ток признаком отсутствия нагрузки;
        нужен powerbank с режимом малых токов или другой источник.
  \item \textbf{Показания напряжения занижены} --- проверьте коэффициент делителя: измерьте батарею мультиметром и поправьте \code{DIVIDER}.
\end{itemize}
\end{trouble}

\begin{summary}
\begin{itemize}[leftmargin=*]
  \item Источник выбирают по ёмкости, безопасности и способу зарядки; для первых опытов --- powerbank или батарейки AA.
  \item Литиевый аккумулятор используют только с платой защиты и модулем зарядки.
  \item Заряд оценивают по напряжению через делитель и АЦП; при разряде устройство засыпает само.
\end{itemize}
\end{summary}

\begin{selfcheck}
\begin{enumerate}
  \item Почему литиевый аккумулятор нельзя использовать без платы защиты?
  \item Зачем перед АЦП ставят делитель и почему его резисторы берут по \SI{100}{k\ohm}, а не по \SI{100}{\ohm}?
  \item Сколько примерно проработает устройство со средним током \SI{20}{mA} от аккумулятора \SI{1000}{mAh}?
\end{enumerate}
\end{selfcheck}

\begin{tasks}
\begin{enumerate}
\item \lv{1} Измерьте напряжение батареи мультиметром и запишите его.
\item \lv{1} Рассчитайте время работы платы от батареи по её ёмкости и измеренному току.
\item \lv{2} Измеряйте напряжение батареи через делитель и АЦП и сравните с мультиметром.
\item \lv{2} Покажите процент заряда на экране по таблице напряжений.
\item \lv{3} Записывайте напряжение батареи раз в минуту в NVS или на SD-карту и постройте график разряда.
\item \lv{4} Сделайте автономный датчик температуры, который просыпается раз в 10 минут, и оцените время его работы от батареи.
\end{enumerate}
\end{tasks}
